% The finite interval lemma is proved in the main text without certificate tables.
\documentclass[11pt,a4paper,reqno]{amsart}
\usepackage[T1]{fontenc}
\usepackage{lmodern}
\usepackage{amsmath,amssymb,amsthm,mathtools,booktabs,array}
\usepackage[margin=.9in]{geometry}
\usepackage[expansion=false]{microtype}
\usepackage{xcolor}
\usepackage[colorlinks=true,linkcolor=blue!45!black,citecolor=blue!45!black,
 urlcolor=blue!55!black,pdftitle={A Sharp Dichotomy for Bounded-Gap Sumsets},
 pdfauthor={Johan Land}]{hyperref}
\setlength{\emergencystretch}{3em}
\newtheorem{maintheorem}{Theorem}
\newtheorem{theorem}{Theorem}[section]
\newtheorem{lemma}[theorem]{Lemma}
\newtheorem{proposition}[theorem]{Proposition}
\newtheorem{corollary}[theorem]{Corollary}
\theoremstyle{definition}
\newtheorem{definition}[theorem]{Definition}
\newtheorem*{problem}{Problem}
\theoremstyle{remark}
\newtheorem{remark}[theorem]{Remark}
\newcommand{\NN}{\mathbb N}
\newcommand{\ZZ}{\mathbb Z}
\newcommand{\dd}{\underline d}
\newcommand{\ssum}{\operatorname{SS}}
\newcommand{\ceil}[1]{\left\lceil #1\right\rceil}
\newcommand{\floor}[1]{\left\lfloor #1\right\rfloor}
\newcommand{\ind}{\mathbf 1}
\begin{document}
\title{A Sharp Dichotomy for Bounded-Gap Sumsets}
\author{Johan Land}
\date{September 2026}
\begin{abstract}
Let $k\ge3$ and $1\le d_1<d_2$. We prove that Erd\H{o}s Problem~1112 has an
affirmative answer exactly when $d_2\ge k+1$; in that range the explicit
lacunarity ratio $d_2+2$ works. Below the threshold we prove more: however
rapidly prescribed successive ratios grow, one sequence meets the $k$-fold
sumset of every admissible bounded-gap sequence. The proof combines a
reciprocal Beatty construction, Kneser's density theorem, and an elementary
proof of the sharp finite subset-sum lemma required at the boundary.
\end{abstract}
\maketitle
\section{Introduction}
For a set $A\subset\NN$, write
\[
 kA:=\{a_1+\cdots+a_k:a_1,\ldots,a_k\in A\},
\]
with repetitions allowed. A sequence $A=\{a_1<a_2<\cdots\}$ is
\emph{$(d_1,d_2)$-admissible} if
$d_1\le a_{n+1}-a_n\le d_2$ for every $n$.
We consider the following formulation recorded as Erd\H{o}s Problem~1112~\cite{E1112},
inspired by a remark of Erd\H{o}s and Graham~\cite{ErGr80}.
\begin{problem}
Let $k\ge3$ and $1\le d_1<d_2$ be integers. Does there exist an integer $r$
such that every increasing sequence $B=\{b_1<b_2<\cdots\}$ of positive
integers satisfying $b_{n+1}\ge rb_n$ admits an increasing sequence
$A=\{a_1<a_2<\cdots\}$ of positive integers with
\[
 d_1\le a_{n+1}-a_n\le d_2\quad(n\ge1),\qquad kA\cap B=\varnothing?
\]
\end{problem}
The following theorem determines exactly when such a ratio exists.

\begin{maintheorem}[Sharp dichotomy]\label{thm:main}
Let $k\ge3$ and $1\le d_1<d_2$ be integers.
\begin{enumerate}
\item If $d_2\ge k+1$, every sequence of positive integers
$B=\{b_1<b_2<\cdots\}$ with $b_{n+1}\ge(d_2+2)b_n$ admits a
$(d_1,d_2)$-admissible $A$ such that $kA\cap B=\varnothing$.
\item If $d_2\le k$, no fixed ratio works. More strongly, for every
sequence $R_1,R_2,\ldots$ of positive integer lower bounds there is one
increasing sequence $B$ of positive integers with $b_{n+1}\ge R_nb_n$
which meets $kA$ for every $(d_1,d_2)$-admissible $A$.
\end{enumerate}
Thus the answer is yes exactly when $d_2\ge k+1$.
\end{maintheorem}
The earlier proof~\cite{Land26}, formally verified in Lean, established this
dichotomy with the bound $192d_2$. The present proof improves that bound and
replaces the finite certificate argument with an elementary construction.

Above the threshold, a reciprocal Beatty sequence avoids a prescribed lacunary
sequence. Below it, every admissible walk has a congruence-class tail in its
$k$-fold sumset; a diagonal construction then meets every such sumset. The
finite interval lemma, Theorem~\ref{thm:sharp}, supplies the required offsets
at the boundary. Its proof is given in Section~\ref{sec:sharp-proof}.
Appendix~\ref{sec:formal} records the correspondence with the Lean formalization.

\medskip\noindent\emph{Acknowledgment.}
I am deeply grateful to Stijn Cambie for his substantial contributions to
simplifying the proof and for providing the manuscript on which this shortened
presentation is based.

\section{The existence construction}
Put $D=d_2$ and assume $D\ge k+1$.
\begin{lemma}[Lacunary interpolation]\label{lem:interpolation}
Let $I=[u,v]\subset(0,1)$ have length $\ell>0$, let $J$ be a
nondegenerate closed bounded real interval, and let positive integers
$q_1<q_2<\cdots$ satisfy $q_{n+1}\ge\rho q_n$. If
\begin{equation}\label{eq:interpolation-cost}
 \ell(\rho-1)\ge1,
\end{equation}
then, after deleting finitely many $q_n$, some $\theta\in J$ satisfies
$\{q_n\theta\}\in I$ for every remaining $n$.
\end{lemma}
\begin{proof}
For a positive integer $q$, the set
\[
 E_q=\{\theta:\{q\theta\}\in I\}
 =\bigcup_{m\in\ZZ}[(m+u)/q,(m+v)/q]
\]
has components of length $\ell/q$, with left endpoints spaced by $1/q$.
For some $N$, the interval
$[q_N\inf J-u,q_N\sup J-v]$ has length $q_N|J|-\ell\ge1$,
so it contains an integer. Thus a component $K_N$ of $E_{q_N}$ lies in $J$.
If $K_n$ is a component of $E_{q_n}$, the allowable indices of components
of $E_{q_{n+1}}$ contained in $K_n$ form a real interval of length
\[
 q_{n+1}|K_n|-\ell
 =\ell(q_{n+1}/q_n-1)\ge1.
\]
It contains an integer, giving $K_{n+1}\subseteq K_n$. The nested nonempty
compact intervals have a common point $\theta$, as required.
\end{proof}
\begin{proposition}[Existence]\label{prop:existence}
If $D\ge k+1$, the ratio $D+2$ works.
\end{proposition}
\begin{proof}
Given $B$ with $b_{n+1}\ge(D+2)b_n$, put
\[
 \varepsilon=\frac1{2D(D+1)},\quad
 I=[\varepsilon,\varepsilon+1/(D+1)],\quad
 J=[1/D,1/D+\varepsilon/(D-1)].
\]
Lemma~\ref{lem:interpolation} applies with equality in
\eqref{eq:interpolation-cost}. It yields $\theta\in J$ and $N_0$ such that
\begin{equation}\label{eq:fractional-safe}
 \{b_n\theta\}\in I\qquad(n\ge N_0).
\end{equation}
Put $\alpha=1/\theta$. Then
\begin{equation}\label{eq:slope-range}
 D-1<\alpha\le D,
\end{equation}
and, since $k\le D-1$,
\begin{equation}\label{eq:room}
 1-k\theta\ge1-(D-1)\max J
 =1/D-\varepsilon=\varepsilon+1/(D+1)=\max I.
\end{equation}
Choose an integer $m\ge1$ with
$k\floor{\alpha m}>\max_{n<N_0}b_n$, reading the empty maximum as zero, and set
\[
 A=\{\floor{\alpha(m+i)}:i=0,1,2,\ldots\}.
\]
Its gaps are $D-1$ or $D$, hence admissible, and its least $k$-fold sum
exceeds every discarded term. If a remaining $b$ were a sum of $k$ such
floors, then for an integer $S$ (the sum of their indices)
$b\le\alpha S<b+k$, whence
\[
 b\theta\le S<(b+k)\theta.
\]
But $0<\{b\theta\}\le1-k\theta$ by \eqref{eq:fractional-safe}--\eqref{eq:room},
so this half-open interval contains no integer. This is a contradiction.
\end{proof}
\begin{remark}
The argument works for every real ratio $\rho>1+D/(D-k)$: take $J$
sufficiently close to $1/D$ and a closed interval
$I\subset(0,1-k\sup J)$ of length $1/(\rho-1)$.
\end{remark}

\section{Tools for tail covering}
Call $A$ \emph{$k$-tail-covering} if, for integers $q\ge1$, $0\le r<q$,
and $X_0$,
\begin{equation}\label{eq:tail-covering}
 \{x\ge X_0:x\equiv r\pmod q\}\subseteq kA.
\end{equation}
\subsection{Density and finite subset sums}
For $X\subseteq\NN_0$ write
$\dd(X)=\liminf_{N\to\infty}|X\cap[1,N]|/N$.
\begin{lemma}[A consequence of Kneser's theorem]\label{lem:density}
Let $P=\{p_0<p_1<\cdots\}$ with $p_0=0$. If $p_n\le cn+C$
eventually, for constants $0<c<k$ and $C$, then $kP$ contains a
congruence-class tail.
\end{lemma}
\begin{proof}
For large $N$, take $n=\floor{(N-C)/c}$. Then $p_n\le N$, so
$\dd(P)\ge1/c>1/k$. The lower-density consequence of Kneser's theorem
\cite{Kneser53,ENS88} needed here is that either $\dd(kP)\ge k\dd(P)$
or $kP$ contains a congruence-class tail. The first alternative is impossible
because density is at most one, proving the claim.
\end{proof}
For a finite multiset $S$, write
$\ssum(S)=\{\sum_{x\in T}x:T\text{ is a submultiset of }S\}$.
\begin{theorem}[Finite interval lemma]\label{thm:sharp}
If $G$ is a finite set of positive integers with $|G|\ge3$, $\gcd(G)=1$,
and $M=\max G$, some multiset of at most $M-1$ elements of $G$ has
$M$ consecutive subset sums.
\end{theorem}
We prove this in Section~\ref{sec:sharp-proof}. First we explain how an interval
of subset sums yields representations of all sufficiently large integers.
\subsection{A slot lemma}
For a gap word $g_1,g_2,\ldots$, put $p_0=0$ and
$p_n=g_1+\cdots+g_n$. A tail of the walk is a translate of $P=\{p_n:n\ge0\}$.
\begin{lemma}[Slots]\label{lem:slots}
Suppose all gaps belong to a finite $G$, every member recurs infinitely
often, $|G|\ge3$, $\gcd(G)=1$, and $M=\max G\le k$. Then $kP$
contains every sufficiently large integer.
\end{lemma}
\begin{proof}
By Theorem~\ref{thm:sharp} choose $S=\{\delta_1,\ldots,\delta_m\}$, $m\le M-1\le k-1$,
with $[c,c+M-1]\subseteq\ssum(S)$. Choose distinct occurrences
$g_{n_j}=\delta_j$. In summand slot $j$ choose $p_{n_j-1}$ or $p_{n_j}$.
These slots realize a fixed base plus every offset in $[c,c+M-1]$.
Set another $k-1-m$ summands equal to $p_0=0$ and let the last summand vary over $p_n$.
For a large target $x$, the residual $x-\mathrm{base}-p_n$ decreases to
$-\infty$ in jumps at most $M$. Its first value at most $c+M-1$ is at
least $c$, so the slots supply the residual. This represents $x$ in $kP$.
\end{proof}

\section{Every subcritical walk is tail-covering}
\begin{proposition}\label{prop:walk}
For $k\ge3$, every increasing sequence of positive integers with gaps at
most $k$ is $k$-tail-covering.
\end{proposition}
\begin{proof}
Discard a finite prefix so that the gaps lie in the nonempty tail alphabet
$G$ of values recurring infinitely often. Put $p_0=0$ and
$p_n=g_1+\cdots+g_n$. Divide all gaps and prefix sums by $g=\gcd(G)$.
A class tail modulo $q$ rescales to one modulo $gq$; translating back only
changes the residue. Relabel so that $\gcd(G)=1$.
If $M=\max G<k$, then $p_n\le Mn$, and Lemma~\ref{lem:density} applies.
Thus assume $M=k$. If $|G|\ge3$, use Lemma~\ref{lem:slots}.
It remains to consider an alphabet with two elements: a singleton would have
gcd $k>1$. Hence $G=\{\delta,k\}$ with $1\le\delta<k$ and $\gcd(\delta,k)=1$.
Set $e=k-\delta$, encode gap $k$ by $h_n=1$ and gap $\delta$ by $h_n=0$,
and write $q_n=h_1+\cdots+h_n$. Then
\begin{equation}\label{eq:prefix-binary}
 p_n=\delta n+e q_n.
\end{equation}
For $s\ge0$ set
\[
 W_s=\left\{\sum_{j=1}^kq_{i_j}:i_j\ge0,\ \sum_{j=1}^ki_j=s\right\}.
\]
This is an integer interval $[w^-(s),w^+(s)]\cap\ZZ$: compositions are
connected by unit transfers, each changing the value by $-1,0,$ or $1$.
A path between extremizers misses no intermediate value. Adding or deleting
one unit shows that each endpoint rises by $0$ or $1$ as $s$ increases by one.
Write $\operatorname{width}(s)=w^+(s)-w^-(s)$.

\medskip\noindent\emph{Sweep.} If $\operatorname{width}(s)\ge k-1$
eventually, then $kP$ contains every large integer. Indeed, a target $x$
is represented at budget $s$ exactly when $w(s)=(x-\delta s)/e\in W_s$.
The integral values of $w(s)$ occur in one class of $s$ modulo $e$.
Along this class $w$ drops by $\delta$ per step, whereas either endpoint
rises by at most $e$. Since $w^+(s)\le s$, for large $x$ the target starts
above the band at the first admissible budget beyond the width threshold;
eventually it becomes negative. If it skips the band, consecutive admissible
budgets satisfy
\[
 w^+(s)<w(s)=w(s+e)+\delta
 \le w^-(s+e)-1+\delta\le w^-(s)+e-1+\delta.
\]
Thus $\operatorname{width}(s)\le k-2$, a contradiction.

\medskip
A binary word is \emph{balanced} if equal-length factors differ in their
numbers of ones by at most one. Define
\[
 V(\sigma)=\max_{i+j=\sigma}(q_i+q_j)-\min_{i+j=\sigma}(q_i+q_j).
\]
If $V(\sigma)\le1$ for every $\sigma$, the word is balanced: comparing
factors with endpoint indices $i,i'$ and $j',j$, where
$i+j=i'+j'$, bounds their count difference by $V(i+j)$.
Let $z_n=n-q_n$ count zeros. Then
\begin{equation}\label{eq:zero-frequency}
 p_n=kn-ez_n.
\end{equation}
If $\liminf z_n/n>0$, choose $0<\gamma\le1$ with $z_n\ge\gamma n$
eventually; then $p_n\le(k-e\gamma)n$, and Lemma~\ref{lem:density} applies.
Assume $\liminf z_n/n=0$. There are infinitely many $\sigma$ with
$V(\sigma)\ge2$. Otherwise a sufficiently late tail would be balanced
by the same endpoint comparison. Since zero recurs, it has a length-$L$
factor containing two zeros. Every length-$L$ factor would then contain a
zero, giving $z_n\ge n/L-O(1)$, a contradiction.

If $k=2u+1$, fix $\sigma_0$ with $V(\sigma_0)\ge2$. For all large $s$,
split the slots into $u$ pairs of subtotal $\sigma_0$ and a singleton
carrying the rest. Independent variation gives width at least
$uV(\sigma_0)\ge k-1$, so the sweep applies.

For even $k$, pairing all the summands leaves no singleton to absorb the
remaining index sum. We instead show that repeated failures of the width
bound would force the binary word to be periodic.
Write $k=2u$, $u\ge2$, and choose $\sigma_0<\sigma_1$ with
$V(\sigma_i)\ge2$, and put $\Delta=\sigma_1-\sigma_0$.
Suppose at a sufficiently large $s$ the width is at most $k-2$, and put
$\tau=s-(u-1)\sigma_0$. Using $u-1$ pairs of subtotal $\sigma_0$ and one
of subtotal $\tau$ gives width at least $k-2+V(\tau)$, so $V(\tau)=0$.
Using $u-2$ pairs of subtotal $\sigma_0$, one of subtotal $\sigma_1$, and
one of subtotal $\tau-\Delta$ likewise gives $V(\tau-\Delta)=0$.
Now $V(n)=0$ implies $q_i+q_{n-i}$ is constant. Successive differences
give $h_{i+1}=h_{n-i}$, so the prefix of length $n$ is a palindrome.
The prefixes of lengths $\tau$ and $\tau-\Delta$ are therefore palindromes;
the shorter is also a suffix of the longer, which consequently has period
$\Delta$. Infinitely many width failures would give such periodic prefixes
of unbounded length, forcing the whole word to have period $\Delta$.
A period contains a zero, contradicting $\liminf z_n/n=0$.
Thus width is eventually at least $k-1$, and the sweep finishes.
\end{proof}

\section{Proof of the non-existence direction}
The preceding proposition gives a congruence-class tail for each admissible
sequence. The following diagonal construction meets all such tails while
allowing arbitrary prescribed growth.
\begin{lemma}\label{lem:certificate}
If every admissible walk is $k$-tail-covering, then any prescribed positive
integer ratios $R_n$ admit one increasing $B$ with $b_{n+1}\ge R_nb_n$
meeting $kA$ for every admissible $A$.
\end{lemma}
\begin{proof}
List pairs $(q,r)$, $q\ge1$, $0\le r<q$, with each recurring infinitely
often, by concatenating the finite lists with $q\le N$ for $N=1,2,\ldots$.
Write the list as $(q_n,r_n)$. Choose $b_1>0$ in its prescribed class and
let $b_{n+1}$ be the least integer in the next class exceeding
$\max\{R_nb_n,b_n,n\}$. The sequence grows as required. For any $A$,
its class in \eqref{eq:tail-covering} occurs infinitely often; the
corresponding $b_n$ are unbounded, so one exceeds $X_0$ and belongs to $kA$.
\end{proof}
\begin{proof}[Proof of Theorem~\ref{thm:main}]
Part (1) is Proposition~\ref{prop:existence}. For $d_2\le k$,
Proposition~\ref{prop:walk} and Lemma~\ref{lem:certificate} give part (2).
\end{proof}

\section{Proof of the finite interval lemma}\label{sec:sharp-proof}
We first reduce Theorem~\ref{thm:sharp} to triples of pairwise coprime
integers. Two constructions then cover the remaining triples, according to
whether their two largest elements sum to a multiple of the smallest.
Write $u^m$ for the multiset of $m$ copies of the positive integer $u$.
\subsection{Reduction to three generators}
\begin{lemma}[Coprime rectangle]\label{lem:rectangle}
If $1\le p<q$ and $\gcd(p,q)=1$, then
\[
 \ssum(p^{q-1},q^{p-1})\supseteq[(p-1)(q-1),pq-1].
\]
The interval has $p+q-1$ terms, at multiset cost $p+q-2$.
\end{lemma}
\begin{proof}
For $n$ in the interval choose $0\le i<q$ with $ip\equiv n\pmod q$.
Then $j=(n-ip)/q$ satisfies $-1<j<p$, so $0\le j\le p-1$.
\end{proof}
\begin{lemma}[Interlacing]\label{lem:interlace}
Let $d,q\ge1$, $\gcd(d,q)=1$, and $q\le N$. If $[A,A+N-1]\subseteq\ssum(C)$,
then scaling $C$ by $d$ and adjoining $d-1$ copies of $q$ gives $N$
consecutive subset sums.
\end{lemma}
\begin{proof}
For $0\le n<N$ choose $0\le j<d$ with
$jq\equiv(d-1)q+n\pmod d$, and put $t=((d-1)q+n-jq)/d$.
Since $0\le(d-1)q+n-jq<dN$, we have $0\le t<N$. Thus
\[
 d(A+t)+jq=dA+(d-1)q+n
\]
supplies the required interval.
\end{proof}
\begin{lemma}[Two-generator targets]\label{lem:long-target}
For coprime $1\le p<h$ and $N\ge2h$, at most $N-\floor{N/h}$ terms
from $\{p,h\}$ have $N$ consecutive subset sums.
\end{lemma}
\begin{proof}
Start with the rectangle's interval of length $p+h-1$. Write $N=sh+r$,
$s\ge2$, $0\le r<h$, and add $t=s-1+\ind_{\{r\ge p\}}$ copies of $h$.
The translates overlap or abut and cover at least $N$ integers, at cost
$h+p+s-3+\ind_{\{r\ge p\}}$. At $s=2$ its slack below $N-s$ is
$h+r-p-1-\ind_{\{r\ge p\}}\ge0$; increasing $s$ increases slack by $h-2$.
\end{proof}
\begin{proposition}[Reduction to three generators]\label{prop:funnel}
Assuming Theorem~\ref{thm:sharp} for maxima below $M$, it suffices at $M$ to treat triples
\begin{equation}\label{eq:funnel}
 G=\{a<b<M\},\quad
 \gcd(a,b)=\gcd(a,M)=\gcd(b,M)=1,\quad a+b\ge M+2.
\end{equation}
\end{proposition}
\begin{proof}
Set $H=G\setminus\{M\}$, $d=\gcd(H)$, and $e=\max(H/d)$.
If $|H|\ge3$, induction gives $e$ consecutive subset sums at cost $e-1$
from $H/d$. Adjoin $\ceil{M/e}-1$ copies of $e$, obtaining at least $M$
consecutive sums. Scale by $d$ and interlace with $d-1$ copies of $M$;
$\gcd(d,M)=1$. The cost is $C=e+d+\ceil{M/e}-3$. Writing $M=de+s$
with $s\ge1$ gives
\[
 M-1-C=(d-1)(e-2)+s-\ceil{s/e}\ge0,
\]
since $e\ge3$.
Otherwise $H=\{a,b\}$, $a<b$. If $d=\gcd(a,b)>1$, apply
Lemma~\ref{lem:long-target} to $\{a/d,b/d\}$ with target $M$, scale,
and interlace with $d-1$ copies of $M$. Here $M>b$, $d\ge2$, and
$\floor{M/(b/d)}\ge d$, so the cost is at most $M-1$.
Now assume $\gcd(a,b)=1$. If $b\le M/2$, the same target lemma costs
at most $M-2$. If $b>M/2$ and $a+b\le M$, one more copy of $b$ extends
the rectangle to length $a+2b-1\ge M$, at cost at most $M-1$.
If $a+b=M+1$, the rectangle itself suffices. Thus $a+b\ge M+2$ remains.
Finally, if $d=\gcd(a,M)>1$, apply the target lemma to $\{a/d,M/d\}$
with target $M$, at cost at most $M-d$, scale, and interlace with $d-1$
copies of $b$; $\gcd(d,b)=1$. The case $\gcd(b,M)>1$ is symmetric.
\end{proof}
For a triple in \eqref{eq:funnel}, put
\begin{equation}\label{eq:triple-notation}
 e=b-a,\qquad h=M-b,\qquad\mu=M-a=e+h.
\end{equation}
Then
\begin{equation}\label{eq:h-range}
 1\le h\le a-2,
\end{equation}
so $a\ge3$ and $\gcd(a,e)=\gcd(a,\mu)=1$.

\subsection{Paired movers}
\begin{lemma}[Paired movers]\label{lem:movers}
If $3\le a<b<M$, $\gcd(a,b)=1$, and $b+M=ra$, then for
\[
 x=r-1,\qquad y=\ceil{(M-r)/2},\qquad z=\floor{(M-r)/2},
\]
$\ssum(a^x,b^y,M^z)$ contains $M$ consecutive integers and $x+y+z=M-1$.
\end{lemma}
\begin{proof}
Here $r\ge3$, and
\begin{equation}\label{eq:mover-bounds}
 M\ge a+r-1,\qquad M\le a(r-1)-1:
\end{equation}
indeed $2M\ge ra+1\ge2a+2r-2$ and $M=ra-b$.
Put $p=\floor{a/2}$, $q=\floor{(a-1)/2}$, and $L=\max(pb,qM)$.
Then $p+q=a-1$, $p\le y$, and $q\le z$.
For $n\in[L,L+M-1]$ choose $w\in[0,a-1]$ with $wb\equiv n\pmod a$.
If $w\le p$, use $w$ unpaired copies of $b$ and divide
$(n-wb)/a=i+r\kappa$, $0\le i<r$. Otherwise put $w'=a-w\le q$,
use $w'$ copies of $M$, and divide $(n-w'M)/a=i+r\kappa$.
The numerators are nonnegative, so $\kappa\ge0$. Each of $\kappa$ further
pairs $(b,M)$ contributes $ra$, and $i\le x$.
The copies are available provided
\begin{align}
 L+M-1+pM&\le a(ry+x),\label{eq:mover-y}\\
 L+M-1+qb&\le a(rz+x).\label{eq:mover-z}
\end{align}
For example, in the first branch
$a(i+r\kappa)=n-wb$ and $a(i+r(\kappa+w))=n+wM$;
these inequalities bound $\kappa\le z$ and $\kappa+w\le y$.
The second branch is symmetric.
If $L=qM$, the inequalities reduce to
$aM-1\le a(ry+x)$ and $qra+M-1\le a(rz+x)$.
The latter uses $q\le z$ and $M-1\le ax$; the former uses
\[
 ry+x\ge r(M-r)/2+r-1\ge M,
\]
since $M\ge a+r-1\ge r+2$.
If $L=pb$, the first inequality follows from $p\le y$ and $M-1\le ax$.
The second is $(a-1)b+M-1\le a(rz+x)$. Odd $a$ gives $p=q$ and
$qM>pb$, so take $a=2v$, $v\ge2$. Put $\tau=M-vr=vr-b\ge1$.
The slack is
\[
 D=2vrz-2(v-1)b-2v+1
  =2vr(z-v+1)+2(v-1)\tau-2v+1.
\]
Here $z=\floor{(r(v-1)+\tau)/2}\ge v-1$. If $z\ge v$, then
$D\ge2vr-1>0$. If $z=v-1$, then $\tau\ne1$ (that would give $z\ge v$),
so $\tau\ge2$ and $D\ge2v-3>0$. Finally $x+y+z=M-1$.
\end{proof}

\subsection{The symmetric residue path}
\begin{lemma}[Centered frame]\label{lem:frame}
If $C$ has a subset sum in one common interval $[L,U]$ for each residue
modulo $a$, put $x=\ceil{(M-1+U-L)/a}$. Then
$\ssum(C\cup a^x)$ contains $[U,L+ax]$, with at least $M$ integers.
\end{lemma}
\begin{proof}
Subtract the chosen representative of $n\pmod a$ from
$U\le n\le L+ax$. The remainder is a nonnegative multiple of $a$ at most $ax$.
\end{proof}
\begin{lemma}[Symmetric $\eta$ argument]\label{lem:eta}
If $a<b<M$ satisfy \eqref{eq:h-range}, are pairwise coprime, and
$a\nmid b+M$, at most $M-1$ terms from $\{a,b,M\}$ have $M$ consecutive
subset sums.
\end{lemma}
\begin{proof}
Let $p,q$ be units modulo $a$, let $\eta\in[0,a-1]$ represent $qp^{-1}$,
and put
\begin{equation}\label{eq:path-data}
 t=\min(\eta,a-\eta),\quad Z=\floor{(a-1)/t},\quad
 r=a-1-tZ,\quad K=t-1+Z.
\end{equation}
For $(p,q)=(b,M)$ we have $\eta\equiv\mu e^{-1}\pmod a$.
The values $0,-1,1$ would give respectively $a\mid M$, $a\mid b+M$,
and $a\mid h$, so $2\le t\le\floor{a/2}$.
If $\eta=t$, write $0\le n<a$ as $n=tk+j$. The representatives $jp+kq$
lie in $[0,S]$, where
\begin{equation}\label{eq:path-positive}
 S=\max\{(t-1)p+(Z-1)q,rp+Zq\}.
\end{equation}
If $\eta=a-t$, write the consecutive residues $t-a,\ldots,t-1$ as
$n=j-tk$; representatives lie in $[0,S]$ with
\begin{equation}\label{eq:path-negative}
 S=(t-1)p+Zq.
\end{equation}
All are subset sums of $p^{t-1},q^Z$. The centered frame with $L=0$ costs
\begin{equation}\label{eq:path-cost}
 \mathcal B=K+\ceil{(M-1+S)/a}.
\end{equation}
We can also reverse $(p,q)$, since $M$ is a unit modulo $a$.
Always
\begin{equation}\label{eq:K-bound}
 K\le\ceil{a/2}.
\end{equation}
Indeed, for $c=\ceil{a/2}$ the concave function $t(c-t+2)$ exceeds $a-1$
at both endpoints of $[2,\floor{a/2}]$, so
$\floor{(a-1)/t}\le c-t+1$. Since $S\le KM$, the budget holds whenever
\begin{equation}\label{eq:crude-budget}
 (a-K-1)M\ge a(K+1)-1.
\end{equation}
For $a=3$ all unit residues are already excluded. We now check the bound
$\mathcal B\le M-1$ separately for odd and even $a$. In each case,
\eqref{eq:crude-budget} handles the smaller values of $K$; only its largest
possible value requires a more precise choice of representatives.

\smallskip\noindent\emph{Odd $a=2n+1$.}
If $K\le n$, the slack in \eqref{eq:crude-budget} at $K=n$ is
$n(\mu-2)\ge0$. Furthermore
$t-1+\floor{2n/t}=n+1$ forces $t\in\{2,n\}$, since
$(t-2)(t-n)<0$ for $2<t<n$. Only $\eta\in\{2,-2,n,n+1\}$ remains.
If $\eta=2$, then $e-h=\lambda a$ for $\lambda\ge0$ by \eqref{eq:h-range}.
Here $Z=n$, $r=0$, $S=nM$. With $\epsilon=0$ for $h=1$ and $\epsilon=1$
for $h\ge2$, evaluation gives
\begin{equation}\label{eq:eta-two}
 (M-1)-\mathcal B=n\lambda+h-2-\epsilon.
\end{equation}
This is nonnegative if $\lambda\ge1$, or if $\lambda=0$ and $e=h=g\ge3$.
For $g=1$, take $n+1$ copies of $a$, one of $a+1$, and $n$ of $a+2$.
At selection level $n+1$, offsets fill $[0,a]$; at the next level offsets
$[1,a]$ remain. The intervals abut, of total length $2a+1\ge a+2$.
For $g=2$, take $n+3$ copies of $a$, one of $a+2$, and $n$ of $a+4$.
With $c=n+1$, each level $\ell\in\{c,c+1,c+2\}$ contains
$\ell a+2[0,a]$. On $[(c+1)a,(c+3)a]$, the middle level supplies one
parity and the other levels supply the other. The costs are respectively
$a+1$ and $a+3$, as required.

If $\eta=-2$, reverse $b,M$; the reversed ratio is $n$.
For $n\ge3$, the normalized residues $1,\ldots,a$ have representatives
\[
 jM\ (1\le j\le n-1),\quad b+jM\ (0\le j\le n-1),\quad 2b,\quad M+2b.
\]
These use $M^{n-1},b^2$ and lie in $[b,b+(n-1)M]$.
The centered frame costs $n+1+\ceil{(nM-1)/a}\le M-1$, since the slack is
\[
 a(M-n-2)-(nM-1)=(n+1)\mu-2n\ge2.
\]
For $n=2$, representatives $M,b,M+b,2b,M+2b$ lie in $[b,M+2b]$.
The cost condition is $2e+3h-9\ge0$. Since $h\equiv2e\pmod5$,
pairwise coprimality and $1\le h\le3$ exclude $e=1,2$, so it holds.

For $\eta=n$ or $n+1$ with $n\ge3$, use the forward path in the first
case and reverse in the second, making its ratio $2$. In either case
$K=n+1$ and $S\le(n-1)b+M$. The budget is at most
$a+1+\ceil{((n+1)e+2h-1)/a}$, and its required slack is
\begin{equation}\label{eq:eta-n}
 ne+(a-2)h-2a+1\ge0.
\end{equation}
For $h\ge2$ this is at least $n-3$. For $h=1$, the first orientation gives
$a\mid3e+2$, so $e\ge3$; the second gives $a\mid e+2$, so $e\ge a-2$.
Both imply the inequality.

\smallskip\noindent\emph{Even $a=2n$.}
If $K\le n-1$, \eqref{eq:crude-budget} holds, with slack $n\mu+1$ at
$K=n-1$. For $a\ge12$ and $3\le t\le n-2$,
$t-1+\floor{(2n-1)/t}\le n-1$. Thus maximal $K=n$ requires
$t\in\{2,n-1,n\}$. As $\eta$ is a unit, only even $n$ and
$\eta=n\pm1$ remain.
For $\eta=n-1$, \eqref{eq:path-positive} gives $Z=2$, $r=1$ and
$S=(n-2)b+M$. The slack is
\begin{equation}\label{eq:eta-even-minus}
 ne+2(n-1)h-2n+1\ge n-1.
\end{equation}
For $\eta=n+1$, $h\equiv ne\pmod{2n}$; $e$ is odd, so
\eqref{eq:h-range} forces $h=n$. Now $S=(n-2)b+2M$ and the slack is
\begin{equation}\label{eq:eta-even-plus}
 2n^2-7n+(n-1)e+1\ge2n(n-3)>0.
\end{equation}
For $a=4,6$ there are no unit residues beyond $\pm1$.
For $a=8$, $\eta=3$ gives $S=b+2M$ and slack
$8(M-5)-(3M+b-1)=4e+5h-7>0$.
For $a=8$, $\eta=5$ forces $h=4$; $S=2b+2M$ gives slack
$8(M-5)-(3M+2b-1)=3e+5h-15>0$.
For $a=10$, reverse if necessary to make the ratio $3$.
The height is at most $2b+2M$, and slack is
$10(M-6)-(3M+2b-1)=5e+7h-9>0$.
This proves the required bound in every case.
\end{proof}
\begin{proof}[Proof of Theorem~\ref{thm:sharp}]
Use strong induction on $M$, with the statement vacuous for $M<3$.
Proposition~\ref{prop:funnel} leaves a triple satisfying \eqref{eq:funnel}.
If $a\mid b+M$, use Lemma~\ref{lem:movers}; otherwise use Lemma~\ref{lem:eta}.
\end{proof}

\begin{samepage}
\section*{Use of AI tools}
Generative AI tools contributed to the development of the original proof,
the drafting and revision of the manuscript, and the construction and debugging
of the Lean formalization.\footnote{In the original development, Claude
(Fable 5 and Opus 4.8) contributed to the mathematical arguments; GPT-5.5 and
Gemini 3.1 assisted with advice and review.}
The shortened presentation is based on a manuscript supplied by Stijn Cambie.
GPT-6-Astra and Claude Sonnet 5 assisted with its revision and formalization.
The formalization was checked by Lean's kernel. The author takes responsibility
for the mathematical claims and their presentation.\par
\end{samepage}

\appendix
\section{Formal verification and correspondence}\label{sec:formal}
Both directions, including the density dependency and finite interval lemma,
are checked in Lean. Selected declarations are listed below; paths are relative
to \texttt{lean/Erdos1112Proof/}.
\begin{center}\small
\begin{tabular}{ll}
\toprule
Paper result & Lean declaration\\
\midrule
\multicolumn{2}{l}{\texttt{Existence/Reciprocal.lean}}\\
Lacunary interpolation & \texttt{reciprocal\_interpolation}\\
Existence with ratio $d_2+2$ & \texttt{existence\_bound\_reciprocal}\\
\multicolumn{2}{l}{\texttt{Short/Intervals.lean}}\\
Coprime rectangle & \texttt{coprime\_rectangle}\\
Interlacing & \texttt{interlacing}\\
Two-generator targets & \texttt{two\_generator\_target}\\
Centered frame & \texttt{centered\_frame}\\
\multicolumn{2}{l}{\texttt{Short/Slots.lean}}\\
Slots and final-summand sweep & \texttt{slots\_from\_run}\\
\multicolumn{2}{l}{\texttt{Short/Normalization.lean}}\\
Tail alphabet and gcd normalization & \texttt{normalized\_walk}\\
\multicolumn{2}{l}{\texttt{Short/Binary.lean}}\\
Zero-liminf binary width argument & \texttt{binary\_eventual\_width}\\
Binary covering from zero liminf & \texttt{binary\_tail\_covering}\\
\multicolumn{2}{l}{\texttt{Short/Funnel.lean}}\\
Reduction to three generators & \texttt{sharpAt\_of\_funnel}\\
\multicolumn{2}{l}{\texttt{Short/Sharp.lean}}\\
Complete finite interval lemma & \texttt{sharp\_all}\\
\multicolumn{2}{l}{\texttt{Short/KneserWeak.lean}}\\
Pairwise density/AP-tail alternative & \texttt{weak\_kneser}\\
\multicolumn{2}{l}{\texttt{Short/Main.lean}}\\
Consequence of Kneser's theorem & \texttt{density\_shortcut}\\
Subcritical walk covering & \texttt{all\_tailCovering}\\
Universal lacunary construction & \texttt{strong\_nonexistence}\\
\bottomrule
\end{tabular}
\end{center}
% lean-correspondence: reciprocal_interpolation | Erdos1112Proof/Existence/Reciprocal.lean
% lean-correspondence: existence_bound_reciprocal | Erdos1112Proof/Existence/Reciprocal.lean
% lean-correspondence: coprime_rectangle | Erdos1112Proof/Short/Intervals.lean
% lean-correspondence: interlacing | Erdos1112Proof/Short/Intervals.lean
% lean-correspondence: two_generator_target | Erdos1112Proof/Short/Intervals.lean
% lean-correspondence: centered_frame | Erdos1112Proof/Short/Intervals.lean
% lean-correspondence: slots_from_run | Erdos1112Proof/Short/Slots.lean
% lean-correspondence: normalized_walk | Erdos1112Proof/Short/Normalization.lean
% lean-correspondence: binary_eventual_width | Erdos1112Proof/Short/Binary.lean
% lean-correspondence: binary_tail_covering | Erdos1112Proof/Short/Binary.lean
% lean-correspondence: sharpAt_of_funnel | Erdos1112Proof/Short/Funnel.lean
% lean-correspondence: sharp_all | Erdos1112Proof/Short/Sharp.lean
% lean-correspondence: weak_kneser | Erdos1112Proof/Short/KneserWeak.lean
% lean-correspondence: density_shortcut | Erdos1112Proof/Short/Main.lean
% lean-correspondence: all_tailCovering | Erdos1112Proof/Short/Main.lean
% lean-correspondence: strong_nonexistence | Erdos1112Proof/Short/Main.lean
The canonical theorems in \texttt{Final.lean} use these proofs throughout.
The Kneser consequence is proved using e-transforms, residue compression and
finite Mann estimates; it is not assumed. The old proofs and certificate tables
have been removed. The definitions in \texttt{lean/Erdos1112.lean} allow repeated
summands, phrase gap bounds additively, and prove equivalence of natural and
integer ratio formulations. The repository supplies build and axiom-audit commands.

\begin{thebibliography}{9}
\bibitem{E1112} T.~F.~Bloom, \emph{Erd\H{o}s Problem \#1112},
\url{https://www.erdosproblems.com/1112}.
\bibitem{ErGr80} P.~Erd\H{o}s and R.~L.~Graham,
\emph{Old and new problems and results in combinatorial number theory},
Monographies de L'Enseignement Math\'ematique \textbf{28},
Universit\'e de Gen\`eve, Geneva, 1980.
\bibitem{ENS88} P.~Erd\H{o}s, M.~B.~Nathanson, and A.~S\'ark\H{o}zy,
\emph{Sumsets containing infinite arithmetic progressions},
J. Number Theory \textbf{28} (1988), 159--166.
\url{https://www.renyi.hu/~p_erdos/1988-33.pdf}
\bibitem{Kneser53} M.~Kneser,
\emph{Absch\"atzung der asymptotischen Dichte von Summenmengen},
Math. Z. \textbf{58} (1953), 459--484.
\url{https://doi.org/10.1007/BF01174162}
\bibitem{Land26} J.~Land,
\emph{A dichotomy for $k$-fold sumsets of bounded-gap sequences
(Erd\H{o}s Problem 1112)}, preprint, 2026.
\url{https://doi.org/10.5281/zenodo.21568276}; source and Lean development:
\url{https://github.com/beetree/math_erdos_1112}.
\end{thebibliography}
\end{document}
